\documentclass{article}
\usepackage{amsmath, amssymb}

\title{Sammanfattning av LWE och modulär linjär algebra}
\author{}
\date{}

\begin{document}

\maketitle

\section{Intuition: varför LWE fungerar}

Bra — "varför" är exakt den viktiga delen här. Det här är själva kärnan i varför LWE-kryptering fungerar.

Vi tar det lugnt och bygger intuition.

\section{Grundidén i LWE-krypto}

Du har:
\[
b = As + e \pmod q
\]

och vill skicka en bit $m \in \{0,1\}$ utan att någon kan lära sig $s$.

\section{Varför använder man slumpvektor $r$?}

När klienten krypterar väljer den en slumpvektor $r$.

Detta gör två viktiga saker:

\subsection*{1. “Mixar” hela systemet}
Den kombinerar många kolumner i $A$ och $b$.

\subsection*{2. Gör varje ciphertext unik}
Samma meddelande ger olika resultat varje gång.

\section{Varför $u = A^T r$?}

\[
u = A^T r
\]

Detta är en slumpmässig linjärkombination av kolumnerna i $A$.

Intuition:
\begin{itemize}
\item $A$ är som ett system av basvektorer
\item $r$ väljer hur du blandar dem
\end{itemize}

Alltså blir $u$ ett "fingeravtryck" av $A$ utan att avslöja något om $s$.

\section{Varför använder man $b^T r$?}

Vi har:
\[
b = As + e
\]

Så:
\[
b^T r = (As + e)^T r
\]

\[
= s^T A^T r + e^T r
\]

Eftersom:
\[
A^T r = u
\]

får vi:
\[
b^T r = s^T u + e^T r
\]

\section{Viktig poäng}

\[
b^T r = s^T u + \text{brus}
\]

Det betyder att:

Den som inte känner $s$ kan inte ta bort denna struktur.

\section{Varför lägger vi till $m \cdot q/2$?}

Vi kodar biten som:
\[
m =
\begin{cases}
0 \rightarrow 0 \\
1 \rightarrow q/2
\end{cases}
\]

Intuition:
\begin{itemize}
\item vi skapar två kluster: nära 0 och nära $q/2$
\end{itemize}

\section{Vad är $v$?}

\[
v = b^T r + m \cdot \frac{q}{2}
\]

Alltså:
\begin{itemize}
\item ett dolt brusigt värde
\item plus en stor skiftning beroende på biten
\end{itemize}

\section{Varför kan mottagaren dekryptera?}

Mottagaren beräknar:
\[
v - s^T u
\]

Sätt in uttrycken:
\[
(s^T u + e^T r + m \cdot q/2) - s^T u
\]

Förenkla:
\[
= e^T r + m \cdot \frac{q}{2}
\]

\section{Resultat}

\[
m \cdot \frac{q}{2} + \text{litet brus}
\]

\section{Beslutsregel}

\begin{itemize}
\item nära $0 \Rightarrow m = 0$
\item nära $q/2 \Rightarrow m = 1$
\end{itemize}

\section{Klassisk LWE-definition}

\[
b = As + e \pmod q
\]

\begin{itemize}
\item $A \in \mathbb{Z}_q^{n \times n}$
\item $s \in \mathbb{Z}_q^n$
\item $e \in \mathbb{Z}_q^n$
\item $b \in \mathbb{Z}_q^n$
\end{itemize}

\section{Intuition (sammanfattning)}

\begin{itemize}
\item $As$ = strukturerad signal
\item $e$ = brus
\item $b$ = förvrängd observation
\end{itemize}

\section{Varför det fungerar}

Sambandet:
\[
b^T r = s^T u + e^T r
\]

gör att hemligheten bara kan elimineras om $s$ är känd.

\section{Slutsats}

LWE bygger på:
\[
b = As + e \pmod q
\]

Säkerheten kommer från svårigheten att lösa linjära system med brus.

\end{document}